\documentclass[11pt]{article}
\usepackage{amsmath, amssymb, amsthm}
\usepackage[margin=1in]{geometry}
\usepackage{algorithm, algpseudocode}

\newcommand{\Gam}{\mathrm{Gam}}
\newcommand{\Poi}{\mathrm{Poi}}
\newcommand{\Cat}{\mathrm{Cat}}
\newcommand{\KL}{\mathrm{KL}}
\newcommand{\E}{\mathbb{E}}
\newcommand{\gb}{\bar{\gamma}}

\title{Poisson SuSiE for One Count Vector}
\author{}
\date{}

\begin{document}
\maketitle

\section{Model}

We observe one count vector $y \in \mathbb{N}^M$ and a known nonnegative
dictionary $F \in \mathbb{R}_+^{K \times M}$. The Poisson SuSiE model writes the
Poisson rate as a sparse sum of $D$ effects:
\[
  y_m \sim \Poi\left(\sum_{d=1}^D \beta_d F_{\gamma_d m}\right),
  \qquad m = 1,\ldots,M.
\]
For each effect $d$,
\begin{align}
  \gamma_d &\sim \Cat\left(\frac{1}{K},\ldots,\frac{1}{K}\right), \\
  \beta_d \mid \gamma_d &\sim \Gam(\alpha_d^0,\lambda_d^0),
\end{align}
where the Gamma distribution is parameterized by shape and rate. The index
$\gamma_d \in \{1,\ldots,K\}$ selects one row of $F$, and $\beta_d > 0$ is the
effect size.

\section{Variational Family}

The approximation factorizes over effects:
\[
  q(\beta,\gamma)
  = \prod_{d=1}^D q(\gamma_d) q(\beta_d \mid \gamma_d),
\]
with
\[
  q(\gamma_d = k) = \gb_{dk},
  \qquad
  q(\beta_d \mid \gamma_d = k) = \Gam(\alpha_{dk},\lambda_{dk}).
\]
We also introduce auxiliary responsibilities $\xi_{dm} \geq 0$ satisfying
$\sum_{d=1}^D \xi_{dm}=1$ for each coordinate $m$. They allocate the count
$y_m$ across the $D$ effects.

The moments used by the updates are
\begin{align}
  \E[\beta_d \mid \gamma_d = k] &= \frac{\alpha_{dk}}{\lambda_{dk}}, \\
  \E[\log \beta_d \mid \gamma_d = k] &= \psi(\alpha_{dk}) - \log \lambda_{dk}, \\
  \E[\log \beta_d] &=
    \sum_{k=1}^K \gb_{dk}\{\psi(\alpha_{dk}) - \log \lambda_{dk}\}, \\
  \E[\log F_{\gamma_d m}] &= \sum_{k=1}^K \gb_{dk}\log F_{km}.
\end{align}

\section{ELBO}

Using Jensen's inequality with the responsibilities $\xi$, the evidence lower
bound is
\begin{align}
  \mathcal{L}(q)
  &= \sum_{d=1}^D \sum_{m=1}^M
    \xi_{dm} y_m
    \left\{\E[\log \beta_d] + \E[\log F_{\gamma_d m}]
    - \log \xi_{dm}\right\} \nonumber \\
  &\quad
    - \sum_{d=1}^D \sum_{k=1}^K
    \gb_{dk}\frac{\alpha_{dk}}{\lambda_{dk}}\sum_{m=1}^M F_{km}
    - \sum_{d=1}^D \KL(q(\gamma_d,\beta_d)\,\|\,p(\gamma_d,\beta_d)).
\end{align}
The KL term decomposes into discrete and continuous parts:
\begin{align}
  \KL_{\mathrm{disc}}
  &= \sum_{d=1}^D \sum_{k=1}^K \gb_{dk}
     \left(\log \gb_{dk} - \log \frac{1}{K}\right), \\
  \KL_{\mathrm{cont}}
  &= \sum_{d=1}^D \sum_{k=1}^K \gb_{dk}
  \Bigg[
    (\alpha_{dk}-\alpha_d^0)\psi(\alpha_{dk})
    + \alpha_d^0(\log \lambda_{dk}-\log \lambda_d^0) \nonumber \\
  &\qquad\qquad
    - \log\frac{\Gamma(\alpha_{dk})}{\Gamma(\alpha_d^0)}
    - (\lambda_{dk}-\lambda_d^0)\frac{\alpha_{dk}}{\lambda_{dk}}
  \Bigg].
\end{align}

\section{Coordinate Updates}

\paragraph{Responsibility update.}
For each coordinate $m$,
\[
  \xi_{dm}
  \propto
  \exp\left\{
    \E[\log \beta_d] + \sum_{k=1}^K \gb_{dk}\log F_{km}
  \right\},
  \qquad d = 1,\ldots,D,
\]
normalized so that $\sum_d \xi_{dm}=1$.

\paragraph{Gamma variational parameters.}
For effect $d$ and factor $k$,
\begin{align}
  \alpha_{dk} &\leftarrow \alpha_d^0 + \sum_{m=1}^M \xi_{dm}y_m, \\
  \lambda_{dk} &\leftarrow \lambda_d^0 + \sum_{m=1}^M F_{km}.
\end{align}
The shape update does not depend on $k$ because the same allocated count is used
for every candidate factor.

\paragraph{Categorical assignment update.}
After dropping terms constant in $k$,
\[
  \log \gb_{dk}
  \leftarrow
  \sum_{m=1}^M \xi_{dm}y_m \log F_{km}
  - \left(\alpha_d^0 + \sum_{m=1}^M \xi_{dm}y_m\right)
    \log\left(\lambda_d^0+\sum_{m=1}^M F_{km}\right),
\]
then $\gb_{d1},\ldots,\gb_{dK}$ are normalized by a softmax.

\section{Empirical Bayes Hyperparameter Update}

When \texttt{update\_prior = TRUE}, the implementation updates
$(\alpha_d^0,\lambda_d^0)$ for each effect. Define
\begin{align}
  \bar{\beta}_d
  &= \sum_{k=1}^K \gb_{dk}\frac{\alpha_{dk}}{\lambda_{dk}}, \\
  \overline{\log \beta}_d
  &= \sum_{k=1}^K \gb_{dk}
     \{\psi(\alpha_{dk})-\log \lambda_{dk}\}.
\end{align}
Maximizing the expected Gamma prior gives
\[
  \lambda_d^0 = \frac{\alpha_d^0}{\bar{\beta}_d},
  \qquad
  \psi(\alpha_d^0) = \overline{\log \beta}_d + \log \lambda_d^0.
\]
The second equation is solved numerically by root finding in the current R
implementation.

\section{Algorithm}

\begin{algorithm}[H]
\caption{Poisson SuSiE}
\begin{algorithmic}[1]
\State Initialize $\gb_{dk}$, $\alpha_{dk}$, and $\lambda_{dk}$.
\For{$t = 1,\ldots,T$}
  \If{symmetry breaking is enabled}
    \State Reset duplicate active effects to the uniform assignment.
  \EndIf
  \For{$d = 1,\ldots,D$}
    \State Update $\xi_{dm}$ for all $d$ and $m$ using the current posterior moments.
    \State Update $\alpha_{dk}$ and $\lambda_{dk}$ for the selected effect $d$.
    \State Update $\gb_{dk}$ for the selected effect $d$ by softmax normalization.
    \If{empirical Bayes is enabled}
      \State Update $\alpha_d^0$ and $\lambda_d^0$.
    \EndIf
  \EndFor
  \State Evaluate the ELBO.
\EndFor
\end{algorithmic}
\end{algorithm}

\section{Effect of Symmetry Breaking}

When $D$ is larger than the number of truly active components, multiple effects
can converge to nearly identical posterior assignment vectors. The optional
symmetry-breaking step detects such duplicate active effects, averages one copy
with the duplicate, and resets the other copy to the uniform assignment over
factors. In the simulation in the workflowr analysis, this keeps the recovered
active factors unchanged up to permutation, while making unused effects more
diffuse instead of duplicating an already-explained factor. The tradeoff is that
the reset is an explicit perturbation of coordinate ascent, so the ELBO can show
small jumps around symmetry-breaking updates.

\end{document}
